\chapter{A finite-resolution certificate for local diffusion geometry}
\label{ch:og-finite}

The ideal conditional theorem involves time derivatives.  A laboratory record
has a nonzero time step, a detector and a finite sample.  This chapter gives a
conservative certificate that pays for all three, under explicit calibration
assumptions.  It uses a single bounded score per stationary three-time record;
no conditional regression, empirical squaring or division by estimated bin
occupancies is required.

\section{Physical and calibration contract}
Let $X$ be a stationary smooth ordinary-reversal diffusion in the model class
of the preceding population theorem. Write $L=S+A$, let $L^\dagger=S-A$ be its
stationary adjoint, and set $f(X_t)=Y_t$.  Assume that $f$ belongs to the domains
needed below and that the Markov semigroups are contractions in the supremum
norm.  The following are calibrated bounds, not quantities silently inferred
from the test records:
\begin{align}
 \|Lf\|_\infty&\le M_{1,+},&
 \|L^\dagger f\|_\infty&\le M_{1,-},\\
 \|L^2f\|_\infty&\le M_{2,+},&
 \|(L^\dagger)^2f\|_\infty&\le M_{2,-}.
\end{align}
Fix a time step $h>0$. A stationary triple is acquired at times $0,h,2h$,
with the middle time relabelled zero. The complete central state determines an
accessible coordinate $U$, a retention indicator $r(U)\in\{0,1\}$, a real
number $q(U)$ and an envelope $a(U)$. On retained states require
\begin{equation}
 |q(U)|\le Q,\qquad a(U)\ge a_{\min}>0,
 \qquad \Gamma_f(X_0)\le a(U) \quad\text{pointwise}.
 \label{eq:og-finite-contract}
\end{equation}
The positive lower bound is on the \emph{chosen upper envelope}; it is not an
assumption that the true local diffusivity is bounded away from zero. The
retention and weighting rules must be fixed before the test records or fitted
on a genuinely independent pilot. A finite partition with constant $q$ and
$a$ on each retained cell is one case. A continuous rule, such as
$q(Y)=-\gamma Y$, is equally admissible.

\section{Uniform control of the time-step bias}
Define the conditional population quantities
\begin{align}
 g_h(U)&=\mathbb E[(Y_h+Y_{-h}-2Y_0)/(2h)\mid U],\\
 K_h(U)&=\mathbb E[(Y_h-Y_0)(Y_{-h}-Y_0)/h^2\mid U].
\end{align}
Write $g(U)=\mathbb E[Sf\mid U]$ and
$K(U)=\mathbb E[(Sf)^2-(Af)^2\mid U]$. Then, almost surely,
\begin{align}
 |g_h-g|&\le\beta_g:=\frac h4(M_{2,+}+M_{2,-}),\\
 |K_h-K|&\le\beta_K:=
 \frac h2(M_{1,-}M_{2,+}+M_{1,+}M_{2,-}).
 \label{eq:og-finite-bias}
\end{align}
For completeness, let
$a_+=(P_hf-f)/h$ and $a_-=(P_h^\dagger f-f)/h$.
The integral semigroup identity gives
\begin{equation}
 \|a_\pm\|_\infty\le M_{1,\pm},\qquad
 \|a_+-Lf\|_\infty\le hM_{2,+}/2,\qquad
 \|a_--L^\dagger f\|_\infty\le hM_{2,-}/2.
\end{equation}
Conditional independence of the stationary past and future given $X_0$
identifies the two-sided product with $a_+a_-$. Apply the exact identity
\begin{equation}
 ab-cd=\tfrac12(a-c)(b+d)+\tfrac12(b-d)(a+c)
\end{equation}
and the preceding norm bounds. This proves the second bound without an
unnecessary term of order $h^2$. Conditioning on $U$ preserves both bounds.

\section{A weighted score with a physical lower-bound guarantee}
For each ideal noiseless triple let
\begin{equation}
 W_h=\frac{r(U)}{a(U)}
 \left[-\left(q(U)-\frac{Y_h-Y_0}{h}\right)
          \left(q(U)-\frac{Y_{-h}-Y_0}{h}\right)
       -2|q(U)|\beta_g-\beta_K\right],
 \label{eq:og-finite-score}
\end{equation}
where excluded centers contribute zero. Then
\begin{equation}
 \mathbb E[W_h]\le\sigma.
 \label{eq:og-finite-main}
\end{equation}
Indeed, the conditional expectation of the uncorrected numerator is
$2qg_h-q^2-K_h$. After the bias penalty it is at most
$2qg-q^2-K\le g^2-K\le\mathbb E[(Af)^2\mid U]$.
The pointwise inequality
$(Af)^2\le\Gamma_f e\le a(U)e$, with
$e=v^\top D^{-1}v$, completes the proof after integration. The integrand
need not be positive; its positive part must not be taken record by record.
Taking the positive part of the final covered lower bound is legitimate
because $\sigma\ge0$.

If $Y$ lies in an interval of width $R$, a common deterministic score range is
\begin{equation}
 B_-=-\frac{(R/h+Q)^2+2Q\beta_g+\beta_K}{a_{\min}},
 \qquad B_+=\frac{R^2}{4h^2a_{\min}}.
 \label{eq:og-finite-range}
\end{equation}
For the upper bound put $u=q-(Y_h-Y_0)/h$ and
$v=q-(Y_{-h}-Y_0)/h$. Then
$-uv\le(u-v)^2/4\le R^2/(4h^2)$. The lower bound uses
$|u|,|v|\le Q+R/h$. The bias penalty is nonnegative and bounded above by the
stated amount. Both bounds include the zero scores from excluded centers.

\section{Finite-sample confidence with exact arithmetic}
Let $W_1,\ldots,W_n$ be independent stationary records under one frozen rule.
Set $B=B_+-B_-$. Hoeffding's inequality gives, with probability at least
$1-\alpha$,
\begin{equation}
 \sigma\ge\left[\overline W-
 B\sqrt{\frac{\log(1/\alpha)}{2n}}\right]_+.
 \label{eq:og-finite-hoeffding}
\end{equation}
A dependence-aware version holds if the conditional mean of each bounded
score, given the preceding complete record, differs from one fixed target
mean by at most a calibrated radius; that radius is subtracted. This is not a
claim that overlapping triples from one trajectory are independent.

For independent identically distributed triples, the empirical-Bernstein
inequality of Maurer and Pontil, Theorem~4 of arXiv:0907.3740v1
\cite{ogMaurer2009}, applied to the
reflected affine-rescaled score, gives for $n\ge2$
\begin{equation}
 \sigma\ge\left[\overline W-
 \sqrt{\frac{2s_n^2\log(2/\alpha)}{n}}
 -\frac{7B\log(2/\alpha)}{3(n-1)}\right]_+,
 \quad s_n^2=\frac1{n-1}\sum_i(W_i-\overline W)^2.
 \label{eq:og-finite-empirical}
\end{equation}
This is a proved empirical-variance inequality, not the substitution of a
sample variance into an inequality requiring the unknown population variance.
The choice of concentration rule is predeclared. Selecting the largest among
several data-dependent lower bounds requires simultaneous coverage or a new
argument; the implementation does not silently grant it.

The reference implementation accepts exact rational score sums and sums of
squares. Integer quantizer codes therefore suffice even for a large record
bank. Suppose each stored score $\widetilde W_i$ differs from the ideal score
by at most $\delta_W$. Apply the selected concentration inequality to the
stored scores in $[B_--\delta_W,B_++\delta_W]$, then subtract
$\delta_W$ from its lower bound. Numerical errors in the empirical mean are
separate penalties; the empirical-Bernstein path requires exact stored means
and second moments. Logarithms and square roots are enclosed outward using
rational arithmetic. An external calibration failure probability may be added to $\alpha$ only
when the physical and concentration arguments hold conditional on frozen
pre-test calibration information, for example from an independent calibration
experiment, or when a separately proved joint guarantee supplies the needed
coverage. A calibration event selected from the same test data does not
automatically preserve independence or coverage. Calibration failure is not
assigned zero merely because a numerical routine runs.

\section{Optional variance calibration}
For noiseless observations, a known global bound $\Gamma_f\le G$ gives a
second concentration route. Conditional It\^o isometry and the bounded drift
imply
\begin{equation}
 \|Y_{\pm h}-Y_0\|_{L^2(\cdot\mid X_0)}
 \le hM_{1,\pm}+\sqrt{2Gh}.
\end{equation}
The two factors in the raw score are conditionally independent given $X_0$.
Consequently the second moment of the corrected score is bounded by
\begin{equation}
 V_*=
 \left[\frac{(Q+M_{1,+}+\sqrt{2G/h})
                   (Q+M_{1,-}+\sqrt{2G/h})
                   +2Q\beta_g+\beta_K}{a_{\min}}\right]^2.
\end{equation}
It can be intersected with the squared maximum absolute range bound.
The usual known-variance Bernstein radius is then
\begin{equation}
 \frac{B\ell}{3n}+
 \sqrt{\frac{2V_*\ell}{n}+\left(\frac{B\ell}{3n}\right)^2},
 \qquad\ell=\log(1/\alpha).
\end{equation}
For fixed calibration bounds, $V_*=O(h^{-2})$ whereas the squared range is
$O(h^{-4})$. This separates a genuine small-step sampling cost from an overly
pessimistic bounded-range estimate. It does not promise useful constants for
an arbitrary apparatus. Bounded score perturbations replace $V_*$ by
$(\sqrt{V_*}+\delta_W)^2$. This option is not used for the noisy-detector
contract below.

\section{Two distinct detector contracts}
First, bounded conditionally unbiased readouts may replace $Y_t$ if the noises
at the three endpoints are conditionally independent given the complete
hidden triple (or satisfy the corresponding zero cross-covariances), and if
$r,q,a$ are measured independently of those noises as functions of the true
central state. The mean drift is unchanged. The product acquires the
additional conditional center variance:
\begin{equation}
 \mathbb E[(Z_h-Z_0)(Z_{-h}-Z_0)\mid X_0]
 =\mathbb E[(Y_h-Y_0)(Y_{-h}-Y_0)\mid X_0]
   +\operatorname{Var}(Z_0\mid X_0).
\end{equation}
Its minus sign in the score makes this noise conservative. Use the detector's
full bounded range in the concentration inequality. This argument does not
justify selecting the rule from the same noisy central reading or estimating
local quadratic variation from a noisy record.

Second, deterministic bounded readout error can be paid as a score
perturbation. Suppose ideal and stored readings are in the same interval and
differ by at most $\delta_Y$. The ideal and approximate scores must use
\emph{identical} $r,q,a$. The factorization of the score gives
\begin{equation}
 |\widetilde W_h-W_h|
 \le\frac{4(R/h+Q)\delta_Y}{h a_{\min}}.
 \label{eq:og-finite-reading-error}
\end{equation}
A useful circle example defines $z_0$ by deterministic rounding of the true
center, fixes $q=-\gamma z_0$, and uses
\begin{equation}
 a(z_0)=\gamma\left[1-\max(0,|z_0|-\delta_Y)^2\right].
\end{equation}
This bounds the true $\Gamma_f=\gamma(1-Y_0^2)$ on the calibrated error event.
Retention based on $z_0$ is the actual fixed state rule, so there is no change
of selector between ideal and stored scores. The reading error, output score
rounding and arithmetic enclosure errors must all be included in
$\delta_W$. A fitted detector model or an uncalibrated numerical simulator is
not a substitute for these premises.

\section{Resources and limits}
A quota of $n$ independent stationary triples uses $3n$ endpoint readings,
$n$ independent stationary preparations and $2hn$ units of declared evolution
time. Excluded centers still consume all these resources. Additional
coordinate readings are counted separately. Geometry calibration, stationary
preparation, resetting, sensor exposure, independent pilot experiments and
energy expenditure require additional accounting. Equal declared evolution
time is not equal total laboratory time or equal energy.

A nonzero step must retain its bias correction. The nonlinear-observable
counterexamples in the adversarial chapter show that an uncorrected
finite-$h$ local ratio can diverge near a zero of $\Gamma_f$, or remain finite
and substantially exceed the true entropy production even after an invertible
scalar recalibration. The lower envelope floor and its sampling
penalty are substantive access requirements. The certificate also retains
the ordinary hidden-state reversal restriction; it does not turn reversible
inertial oscillations into evidence of dissipated heat.
